\documentclass[11pt,a4paper]{article}

\usepackage[utf8]{inputenc}
\usepackage[T1]{fontenc}
\usepackage[polish]{babel}
\usepackage{lmodern}
\usepackage{microtype}
\usepackage[a4paper,margin=2.5cm]{geometry}
\usepackage{parskip}
\usepackage{graphicx}
\usepackage{array}
\usepackage{booktabs}
\usepackage{tabularx}
\usepackage{enumitem}
\usepackage[font=small,labelfont=bf]{caption}
\usepackage{xcolor}
\usepackage{tikz}
\usetikzlibrary{arrows.meta,positioning}
\usepackage{float}
\usepackage{hyperref}

% Generated by export_data.py; do not edit by hand.
\newcommand{\NInspections}{316\,636}
\newcommand{\NPlaces}{49\,270}
\newcommand{\FirstYear}{2010}
\newcommand{\DataUntil}{21.10.2026}
\newcommand{\MandatoryShare}{46}
\newcommand{\MedianDaily}{76}
\newcommand{\RoutinePerDay}{35}
\newcommand{\RefDay}{20.10.2026}
\newcommand{\DueRef}{7\,481}
\newcommand{\NTrain}{70\,027}
\newcommand{\NValid}{8\,571}
\newcommand{\NTest}{13\,447}
\newcommand{\AucValidGB}{0,638}
\newcommand{\BrierValidGB}{0,171}
\newcommand{\TopValidGB}{42,0}
\newcommand{\AucTestGB}{0,626}
\newcommand{\BrierTestGB}{0,168}
\newcommand{\TopTestGB}{35,0}
\newcommand{\AucValidLR}{0,629}
\newcommand{\BrierValidLR}{0,172}
\newcommand{\TopValidLR}{40,8}
\newcommand{\AucTestLR}{0,626}
\newcommand{\BrierTestLR}{0,169}
\newcommand{\TopTestLR}{36,6}
\newcommand{\AllValid}{23,2}
\newcommand{\AllTest}{22,3}
\newcommand{\BrierTestCal}{0,168}
\newcommand{\TopTestCal}{35}
\newcommand{\AllTestCal}{22}
\newcommand{\TopDecilePred}{41,4}
\newcommand{\TopDecileObs}{35,0}
\newcommand{\LocGain}{0,007}
\newcommand{\LocTestAuc}{0,610}
\newcommand{\LocBurdenMax}{1,94}
\newcommand{\LocBurdenMaxQ}{2}
\newcommand{\BurdenModelMin}{0,79}
\newcommand{\BurdenModelMax}{1,52}
\newcommand{\CalGapMax}{3,4}
\newcommand{\FalseAlarmMin}{60}
\newcommand{\FalseAlarmMax}{67}
\newcommand{\AucQMin}{0,61}
\newcommand{\AucQMax}{0,67}
\newcommand{\BurdenModelMaxQ}{2}
\newcommand{\NFeatures}{11}
\newcommand{\NShownFeatures}{8}
\newcommand{\WRisk}{0,57}
\newcommand{\WVuln}{0,29}
\newcommand{\WOverdue}{0,14}
\newcommand{\CRCrit}{0,00}
\newcommand{\CRGroups}{0,01}
\newcommand{\GChildren}{1,00}
\newcommand{\GHealth}{1,00}
\newcommand{\GRest}{0,39}
\newcommand{\GRetail}{0,19}
\newcommand{\GOther}{0,18}
\newcommand{\AhpRV}{2}
\newcommand{\AhpRO}{4}
\newcommand{\AhpVO}{2}
\newcommand{\DaysSys}{5,0}
\newcommand{\DaysHiSys}{23,7}
\newcommand{\DaysRisk}{6,5}
\newcommand{\DaysHiRisk}{5,8}
\newcommand{\DaysAhp}{4,9}
\newcommand{\DaysHiAhp}{23,7}
\newcommand{\DaysRandom}{0,3}
\newcommand{\AhpLoss}{1,6}
\newcommand{\KmSys}{1,7}
\newcommand{\KmAhp}{3,2}
\newcommand{\KmFair}{3,3}
\newcommand{\KmCity}{3,1}
\newcommand{\SysDaysMin}{2,2}
\newcommand{\SysDaysMax}{6,7}
\newcommand{\HalfCity}{48}
\newcommand{\HalfSys}{60}
\newcommand{\HalfRisk}{63}
\newcommand{\CapFull}{1\,236}
\newcommand{\CapEighty}{688}
\newcommand{\CapShareEighty}{56}
\newcommand{\ReachCity}{53}
\newcommand{\ReachSys}{65}
\newcommand{\LateCity}{14,0}
\newcommand{\LateSys}{9,6}
\newcommand{\BurdenCityMin}{0,90}
\newcommand{\BurdenCityMax}{1,22}
\newcommand{\BurdenAhpMin}{0,83}
\newcommand{\BurdenAhpMax}{1,35}
\newcommand{\BurdenFairMin}{0,93}
\newcommand{\BurdenFairMax}{1,18}
\newcommand{\BurdenSysMin}{0,87}
\newcommand{\BurdenSysMax}{1,29}
\newcommand{\BurdenSysPoorest}{1,01}
\newcommand{\QDaysSysMin}{2,0}
\newcommand{\QDaysSysMax}{11,3}


\newcommand{\coursename}{Systemy wspomagania decyzji}
\newcommand{\projecttitle}{Planowanie kontroli sanitarnych w~Chicago}
\newcommand{\projectsubtitle}{Zadanie 1: Karta problemu decyzyjnego}
\newcommand{\authorone}{Wojciech Jasiński}
\newcommand{\authortwo}{Piotr Iśtok}
\newcommand{\lecturer}{dr inż. Grzegorz Kunikowski}
\newcommand{\projectdate}{Warszawa, 2026}

\setlist{itemsep=0.15em, parsep=0pt, topsep=0.3em, partopsep=0pt}
\setlength{\parskip}{0.5em plus 0.1em}
\captionsetup[table]{skip=4pt}
\setlength{\intextsep}{8pt plus 2pt minus 2pt}
\renewcommand{\arraystretch}{1.2}
\newcolumntype{L}[1]{>{\raggedright\arraybackslash}p{#1}}
\newcolumntype{Y}{>{\raggedright\arraybackslash}X}

\makeatletter
\newcommand{\needspace}[1]{\par\begingroup\setlength{\dimen@}{#1}\vskip\z@\@plus\dimen@\penalty -100\vskip\z@\@plus -\dimen@\vskip\dimen@\penalty 9999\vskip -\dimen@\vskip\z@skip\endgroup}
\makeatother

\definecolor{navy}{HTML}{1F4E79}
\definecolor{accent}{HTML}{C9731E}
\hypersetup{colorlinks, linkcolor=navy, urlcolor=navy, citecolor=navy}

\begin{document}

\begin{titlepage}
  \centering
  \includegraphics[width=0.75\textwidth]{../tech-report/politechnika.png}\par
  \vspace{2.5cm}
  {\large \coursename\par}
  \vspace{1.5cm}
  {\LARGE\bfseries \projecttitle\par}
  \vspace{0.5cm}
  {\Large \projectsubtitle\par}
  \vspace{2.5cm}
  {\large Autorzy\par}
  \vspace{0.3cm}
  {\large \authorone\par}
  \vspace{0.1cm}
  {\large \authortwo\par}
  \vspace{1cm}
  {\large Prowadzący\par}
  \vspace{0.3cm}
  {\large \lecturer\par}
  \vfill
  {\large \projectdate\par}
\end{titlepage}

% =============================================================================================
\section{Decyzja i decydent}

Departament Zdrowia Publicznego Chicago kontroluje lokale gastronomiczne: restauracje, sklepy,
stołówki szkolne, kuchnie szpitali i domów opieki. W otwartych danych miasta z lat
\FirstYear--2026 jest \NInspections{} kontroli \NPlaces{} lokali~\cite{foodinsp}. Prawie połowa
kontroli (\MandatoryShare\% w latach 2025--2026) to wizyty obowiązkowe: po skargach, przy
wydaniu licencji i kontrole ponowne. Na kontrole rutynowe zostaje około \RoutinePerDay{} wizyt
dziennie, a lokali z zaległą kontrolą rutynową jest kilka tysięcy. Każdego dnia trzeba więc
wybrać, które z nich odwiedzić najpierw.

Miasto od 2015 roku szereguje lokale według ryzyka modelem uczenia maszynowego; w teście
z 2014 roku pozwolił on wykrywać naruszenia krytyczne średnio o 7,5 dnia wcześniej~\cite{cityeval}.
Model wskazuje jednak tylko, gdzie ryzyko jest największe. Nie uwzględnia skutków naruszenia
dla klientów, wizyt obowiązkowych, dostępnych inspektorów ani dojazdów, a bez nich nie da się
ułożyć planu dnia.

\textbf{Pytanie decyzyjne:} \emph{które lokale skontrolować danego dnia, kto je odwiedzi
i w jakiej kolejności, aby jak najwcześniej wykryć lokale z poważnymi uchybieniami przy
dostępnym zespole i po wykonaniu wizyt obowiązkowych?}

\textbf{Cel:} jak najwcześniejsze wykrywanie lokali, które nie przechodzą kontroli, zwłaszcza
tych, w których jedzą dzieci i osoby chore, przy wykonaniu wizyt obowiązkowych, rozsądnym czasie
dojazdów i równym traktowaniu dzielnic.

Decydentem jest \textbf{kierownik zespołu inspektorów}: codziennie ustala plan i za niego
odpowiada. System przygotowuje propozycję planu, którą kierownik może zmienić, zanim ją
zatwierdzi. Pozostałe role zestawia tabela~\ref{tab:roles}.

\begin{table}[H]
\centering\small
\caption{Interesariusze decyzji.}
\label{tab:roles}
\begin{tabularx}{\textwidth}{@{}L{3.6cm}L{4.6cm}Y@{}}
\toprule
Interesariusz & Rola & Oczekiwania \\
\midrule
Kierownik zespołu & Decydent: ustala i zatwierdza plan & Plan dobrze wykorzystujący zespół,
  który da się uzasadnić \\
Inspektorzy & Wykonują plan i wpisują wyniki & Realistyczne trasy, równe obciążenie \\
Kierownictwo departamentu & Ustala liczbę etatów i zasady kontroli & Mierzalna poprawa
  ochrony zdrowia \\
Mieszkańcy i klienci lokali & Ponoszą skutki opóźnień & Szybkie wykrywanie zagrożeń \\
Właściciele lokali & Podmioty kontrolowane & Przewidywalne, niedyskryminujące traktowanie \\
\bottomrule
\end{tabularx}
\end{table}

Główne napięcie dotyczy skuteczności i sprawiedliwości: kontrole powinny trafiać tam, gdzie
ryzyko jest największe, ale nie mogą nieproporcjonalnie obciążać niektórych dzielnic.

\textbf{Zakres i ograniczenia.} Decyzja dotyczy planu jednego dnia: wyboru lokali, przydziału
do inspektorów i kolejności wizyt. Nie obejmuje przebiegu kontroli, sankcji ani liczby etatów.
Plan musi uwzględniać:
\begin{itemize}
  \item liczbę inspektorów obecnych danego dnia;
  \item czas pracy inspektora: do 6 wizyt i 8 godzin łącznie z dojazdami;
  \item pierwszeństwo wizyt obowiązkowych;
  \item częstotliwość kontroli zależną od kategorii ryzyka lokalu;
  \item zakaz używania cech dyskryminujących i obowiązek uzasadnienia każdej wizyty;
  \item to, że ostateczną decyzję podejmuje człowiek.
\end{itemize}

\textbf{Charakter problemu.} Decyzja jest operacyjna i powtarzana codziennie, częściowo
ustrukturyzowana i podejmowana w warunkach ryzyka: wynik kontroli nie jest znany przed wizytą,
ale można oszacować jego prawdopodobieństwo. Jest wielokryterialna i wymaga optymalizacji przy
ograniczeniach. Odpowiada jej hybrydowy system wspomagania decyzji: model uczony na danych
historycznych, metoda AHP do wyrażenia preferencji i optymalizacja planu, którego wynik
zatwierdza człowiek.

% =============================================================================================
\section{Uzasadnienie}

\textbf{Wartość dla użytkownika.} Kierownik nie musi ręcznie przeglądać tysięcy zaległych
lokali. Dostaje propozycję planu dla całego zespołu z uzasadnieniem każdej wizyty i mapą tras,
a inspektor widzi swoją trasę na dany dzień.

\textbf{Uzasadnienie biznesowe.} Liczba inspektorów jest stała, więc poprawa może wynikać
tylko z lepszej kolejności wizyt i krótszych dojazdów. Wcześniejsze wykrycie uchybień zmniejsza
ryzyko ognisk zatruć, a z nimi kosztów leczenia i dochodzeń epidemiologicznych. W symulacji na danych z lat 2025--2026
plan systemu wykrywał lokale z wynikiem negatywnym średnio o \DaysSys{} dnia wcześniej niż
kolejność stosowana przez miasto, a placówki dla dzieci i opieki zdrowotnej o \DaysHiSys{} dnia
wcześniej. Wynik dotyczy danych historycznych i wymaga potwierdzenia w praktyce.

\textbf{Inne uzasadnienie.} Ochrona zdrowia publicznego, przejrzystość decyzji urzędu
i równe traktowanie mieszkańców różnych dzielnic.

% =============================================================================================
\section{Warianty decyzji}

Dopuszczalnych planów dnia jest zbyt wiele, by porównywać je pojedynczo. Porównujemy więc
reguły, według których plan powstaje (tabela~\ref{tab:variants}). Przy każdej z nich kierownik
może dodać lub pominąć lokal, zmienić skład zespołu i wagi kryteriów.

\begin{table}[H]
\centering\small
\caption{Porównywane warianty.}
\label{tab:variants}
\begin{tabularx}{\textwidth}{@{}L{4.2cm}Y@{}}
\toprule
Wariant & Opis \\
\midrule
W0 Kolejność miasta & Kolejność, w jakiej miasto faktycznie kontrolowało lokale (stan
  obecny) \\
W1 Kolejność losowa & Punkt kontrolny dla metody oceny \\
W2 Tylko ryzyko & Najpierw lokale o najwyższym przewidywanym ryzyku \\
W3 Ryzyko i skutki & Priorytet łączący ryzyko, wrażliwość klientów i zaległość kontroli \\
W4 Pełny system & W3 z równoważeniem dzielnic i trasami inspektorów; wizyty obowiązkowe
  najpierw \\
\bottomrule
\end{tabularx}
\end{table}

% =============================================================================================
\section{Kryteria i preferencje}

Kryteria działają na dwóch poziomach (tabela~\ref{tab:criteria}). Pierwsze trzy ustalają
kolejność lokali w planie. Pozostałe służą do oceny całego planu i porównania wariantów.

Preferencje decydenta wyrażamy metodą AHP~\cite{saaty}. Kierownik porównuje parami kryteria
wyboru lokalu, a z porównań wynikają ich wagi oraz współczynnik zgodności, który ostrzega
o sprzecznych ocenach. Tą samą metodą ustalamy wagi typów lokali. Domyślnie największą wagę ma
ryzyko, mniejszą wrażliwość klientów, a najmniejszą zaległość; kierownik może je zmienić.

\textbf{Przykład ilustracyjny.} W planie zostało jedno miejsce, a do wyboru są bar szybkiej
obsługi z ryzykiem 0,40 i stołówka szpitalna z ryzykiem 0,25. Sama predykcja wskazałaby bar.
Po uwzględnieniu wrażliwości klientów wyżej trafia stołówka, bo uchybienia zagrażają tam
osłabionym pacjentom.

\begin{table}[H]
\centering\small
\caption{Kryteria, ich kierunek i pomiar.}
\label{tab:criteria}
\begin{tabularx}{\textwidth}{@{}L{4.2cm}L{1.9cm}Y@{}}
\toprule
Kryterium & Kierunek & Pomiar \\
\midrule
\multicolumn{3}{@{}l}{\textit{Wybór lokalu}} \\
Ryzyko wyniku negatywnego & max & Prawdopodobieństwo z modelu predykcyjnego \\
Wrażliwość klientów & max & Waga typu lokalu (np. szkoła, szpital, restauracja) \\
Zaległość kontroli & max & Czas od terminu wynikającego z kategorii ryzyka \\
\midrule
\multicolumn{3}{@{}l}{\textit{Ocena planu}} \\
Szybkość wykrycia & max & O ile dni wcześniej niż w W0 znaleziono lokale z wynikiem
  negatywnym \\
Dojazdy & min & Kilometry dojazdu na wizytę \\
Równe traktowanie dzielnic & min odchylenia & Udział wczesnych wizyt w grupach dzielnic
  według dochodu wobec ich udziału w wynikach negatywnych \\
Wizyty obowiązkowe & warunek & Wykonane wszystkie, które mieszczą się w dniu \\
\bottomrule
\end{tabularx}
\end{table}

% =============================================================================================
\section{Dane i niepewność}

Korzystamy wyłącznie z otwartych danych miasta Chicago. Główne źródło to rejestr kontroli
sanitarnych~\cite{foodinsp}: dla każdej kontroli zawiera lokal, jego typ i kategorię ryzyka,
adres i współrzędne, datę, rodzaj i wynik kontroli oraz opis naruszeń. Uzupełniają go granice
dzielnic~\cite{areas} i dochód ich mieszkańców~\cite{socio}, używany wyłącznie do sprawdzania
równego traktowania. W działającym systemie dochodzą wyniki wpisywane przez inspektorów.

Najważniejsze źródła niepewności i sposób ich uwzględnienia:
\begin{itemize}
  \item \textbf{Wynik kontroli} nie jest znany przed wizytą. Model podaje tylko
        prawdopodobieństwo i ma umiarkowaną trafność (AUC~\AucTestGB), dlatego ryzyko jest
        jednym z kilku kryteriów, a plan zatwierdza człowiek.
  \item \textbf{Zgłoszenia i nieobecności} nie są znane z wyprzedzeniem, więc plan powstaje
        na nowo każdego dnia.
  \item \textbf{Dostępność lokalu}: lokal bywa zamknięty lub zlikwidowany. Nieudana wizyta
        nie zamyka zgłoszenia.
  \item \textbf{Dane}: wyniki znamy tylko dla lokali, które skontrolowano, a czasy dojazdu są
        przybliżone.
\end{itemize}

Koszty błędów są nierówne. Pominięcie lokalu z poważnymi uchybieniami może skończyć się
zatruciem klientów, a zbędna kontrola kosztuje tylko czas inspektora. Plan nie odrzuca więc
lokali progiem klasyfikatora, tylko układa je według priorytetu.

\textbf{Przykład.} Wśród 10\% lokali uznanych przez model za najbardziej ryzykowne wynik
negatywny miało \TopTestCal\% wobec \AllTestCal\% w całej próbie. Model dobrze wskazuje, gdzie
szukać, ale około dwie trzecie wskazanych lokali kontrolę przechodzi.

% =============================================================================================
\needspace{9cm}
\section{Szkic rozwiązania}

System składa się z pięciu modułów, a decyzję podejmuje człowiek (rys.~\ref{fig:arch}).

\begin{figure}[H]
\centering
\begin{tikzpicture}[
  box/.style={draw=navy, fill=navy!6, rounded corners=2pt, align=center, font=\footnotesize,
              minimum height=1.2cm, text width=2.05cm, inner sep=3pt},
  hum/.style={draw=accent, fill=accent!10, rounded corners=2pt, align=center, font=\footnotesize,
              minimum height=1.0cm, text width=2.6cm, inner sep=3pt},
  arr/.style={-{Stealth[length=2mm]}, navy, thick},
  harr/.style={-{Stealth[length=2mm]}, accent, thick},
  lab/.style={font=\scriptsize, text=black!70, align=center},
  node distance=0.62cm]
  \node[box] (data) {\textbf{Dane}\\historia kontroli};
  \node[box, right=of data] (pred) {\textbf{Predykcja}\\ryzyko lokalu};
  \node[box, right=of pred] (pref) {\textbf{Preferencje}\\priorytet (AHP)};
  \node[box, right=of pref] (dec) {\textbf{Decyzja}\\plan dnia\\i trasy};
  \node[box, right=of dec] (db) {\textbf{Baza}\\plany i wyniki};
  \node[hum, below=1.3cm of pref] (mgr) {\textbf{Kierownik}\\zmienia\\i zatwierdza plan};
  \node[hum, below=1.3cm of db] (insp) {\textbf{Inspektorzy}\\wykonują wizyty};
  \node[hum, below=1.3cm of data] (city) {\textbf{Dane miasta}};
  \draw[arr] (city) -- (data);
  \draw[arr] (data) -- (pred);
  \draw[arr] (pred) -- (pref);
  \draw[arr] (pref) -- (dec);
  \draw[arr] (dec) -- (db);
  \draw[harr] (mgr) -- node[lab, left, pos=0.45] {wagi} (pref);
  \draw[harr] (mgr.east) -| node[lab, left, pos=0.75] {zmiany,\\zatwierdzenie} (dec.south);
  \draw[harr] (db.south) ++(-0.4,0) -- node[lab, left] {trasa} ([xshift=-4mm]insp.north);
  \draw[harr] ([xshift=4mm]insp.north) -- node[lab, right] {wynik} ([xshift=0.4cm]db.south);
  \draw[arr, rounded corners=4pt] (db.north) -- ++(0,0.5) -| node[lab, above, pos=0.25]
    {wyniki wizyt uzupełniają historię lokali} (data.north);
\end{tikzpicture}
\caption{Moduły systemu (granatowe), ludzie i źródło danych (pomarańczowe) oraz przepływ
danych.}
\label{fig:arch}
\end{figure}

Jeden dzień pracy systemu przebiega następująco:
\begin{enumerate}
  \item Moduł danych pobiera dane miasta, ujednolica je i dopisuje wyniki wpisane w systemie.
  \item Moduł predykcyjny szacuje ryzyko każdego lokalu z historii sprzed danego dnia.
  \item Moduł preferencji łączy ryzyko z wrażliwością klientów i zaległością w priorytet lokalu.
  \item Moduł decyzyjny układa plan: najpierw wizyty obowiązkowe, potem lokale o najwyższym
        priorytecie, rozdzielone na trasy inspektorów.
  \item Kierownik przegląda plan wraz z uzasadnieniem wizyt, poprawia go i zatwierdza.
  \item Inspektorzy wykonują wizyty i wpisują wyniki, które zasilają kolejne plany.
\end{enumerate}

% =============================================================================================
\section{Założenia przygotowania prototypu}

Prototyp to aplikacja przeglądarkowa z dwiema rolami: kierownika, który przegląda
i zatwierdza plan dnia, oraz inspektora, który widzi swoją trasę i wpisuje wyniki. Działa
w trybie demonstracyjnym: wybiera się dzień z przeszłości, a system układa plan tak, jakby
działał w tamtym czasie.
\begin{itemize}
  \item \textbf{Na prawdziwych danych:} model, priorytety, układanie planu i ocena wariantów.
  \item \textbf{Uproszczone:} czasy wizyt i dojazdów; wagi AHP ustala zespół, nie eksperci.
  \item \textbf{Fikcyjne:} zespół inspektorów i wpisywane wyniki. Zgłoszenia danego dnia to
        kontrole obowiązkowe, które miasto wtedy wykonało.
\end{itemize}

\textbf{Założenia i ich weryfikacja} zestawia tabela~\ref{tab:assumptions}.

\begin{table}[H]
\centering\small
\caption{Założenia i ich krytyczna weryfikacja.}
\label{tab:assumptions}
\begin{tabularx}{\textwidth}{@{}L{4.6cm}Y@{}}
\toprule
Założenie & Weryfikacja i ryzyko \\
\midrule
Model jest wystarczająco trafny & AUC \AucTestGB{} na danych z lat 2025--2026, nieużywanych do
  uczenia; trafność umiarkowana, więc model tylko wspiera decyzję. \\
Symulacja jest uczciwą oceną & Kolejność losowa daje wynik bliski zera; symulacja obejmuje
  jednak tylko kontrole, które się odbyły. \\
Model nie dyskryminuje dzielnic & Sprawdzone w grupach dzielnic według dochodu; współrzędne
  lokali usunięto z modelu. \\
Częstotliwość kontroli & Okresy pół roku, rok i dwa lata wymagają potwierdzenia w przepisach
  miasta. \\
Czas wizyty i dojazdu & 45 minut na wizytę i dojazd po siatce ulic to przybliżenia. \\
\bottomrule
\end{tabularx}
\end{table}

\textbf{Organizacja pracy.} Piotr Iśtok: moduł predykcyjny, ocena wariantów i analiza
równego traktowania dzielnic. Wojciech Jasiński: moduł decyzyjny, architektura i aplikacja.

% =============================================================================================
\section{Wykorzystanie GenAI i odpowiedzialność zespołu}

% DO SPRAWDZENIA PRZEZ ZESPÓŁ przed oddaniem.
Przy przygotowaniu tej karty i wcześniejszej prezentacji projektu korzystaliśmy z asystenta
opartego na dużym modelu językowym (Claude, Anthropic). Pomagał w redakcji tekstu,
porządkowaniu struktury dokumentu i sprawdzaniu jego spójności. Wszystkie treści przejrzeliśmy,
a podane liczby sprawdziliśmy w danych źródłowych.

Samodzielnie określiliśmy temat i zakres projektu, problem decyzyjny i decydenta, warianty,
kryteria oraz założenia. Za treść dokumentu odpowiada zespół.

\begingroup\small\raggedright
\begin{thebibliography}{9}
\bibitem{foodinsp} City of Chicago, \textit{Food Inspections}, Chicago Data Portal,
  \url{https://data.cityofchicago.org/d/4ijn-s7e5}.
\bibitem{areas} City of Chicago, \textit{Boundaries -- Community Areas},
  \url{https://data.cityofchicago.org/d/igwz-8jzy}.
\bibitem{socio} City of Chicago, \textit{Census Data -- Selected socioeconomic indicators,
  2008--2012}, \url{https://data.cityofchicago.org/d/kn9c-c2s2}.
\bibitem{cityeval} City of Chicago, \textit{Food Inspections Evaluation},
  \url{https://github.com/Chicago/food-inspections-evaluation}.
\bibitem{saaty} T. L. Saaty, \textit{The Analytic Hierarchy Process}, McGraw-Hill, 1980.
\end{thebibliography}
\endgroup

\end{document}
